\documentclass[11pt]{article}
\usepackage[a4paper,margin=20mm]{geometry}
\usepackage{amsmath,amssymb,amsthm,microtype}
\usepackage[colorlinks=true,linkcolor=blue,urlcolor=blue]{hyperref}
\newtheorem{theorem}{Theorem}
\title{Joint polynomial approximation of MPUs with support at most two}
\author{TNLean contributors}
\date{3 October 2026}
\begin{document}
\maketitle

A fixed physical-entry alphabet gives exact joint polynomial
circuits by recovering discrete coefficient labels
\cite{fixed-alphabet}. For support
at most two per row and column, one can remove that alphabet
assumption and obtain arbitrary operator-norm accuracy with
exact auxiliary erasure. Orthogonality restricts the support to
singleton and two-dimensional unitary blocks. A common rounded
edge predicate gives consistent component addresses; exact
unitary approximations inside those components preserve the
address subspace and permit complete auxiliary erasure.
The finite-precision Choi contractions used below follow
\nolinkurl{docs/audits/2026-10-03_mpu_joint_polynomial_permutations.tex}.
These new results are mathematical proofs and are not asserted
to be formalized in Lean.

\section{Joint polynomial approximation with exact erasure at support two}

Fix the local dimension $d\ge2$.
The sparsity assumption below is on the rows and columns of the
full physical unitary in the computational basis, not on its
local MPO tensors.
Gates are arbitrary one- and two-qudit unitaries.

\begin{theorem}
Let $U$ be a unitary on $N\ge1$ $d$-dimensional sites, with at
most two nonzero entries in every row and every column, and with
operator Schmidt ranks at most $D\ge1$ across consecutive cuts.
For every $0<\varepsilon<1$ there is a circuit $C$ and a unitary
$\widetilde U$ such that
\[
 CJ=J\widetilde U,\qquad
 \|\widetilde U-U\|_{\rm op}\le\varepsilon,
\]
where $J$ initializes every auxiliary register at zero.
The gate count and auxiliary space are polynomial jointly in
$N,D,\log(1/\varepsilon)$, including after routing to adjacent
two-qudit gates on a line. Auxiliary erasure is exact.
There is no lower bound on nonzero entries or nonzero Schmidt
values.

This is a nonuniform arbitrary-gate existence statement.
The proof is mathematical and not Lean formalized.
The auxiliary bound is polynomial, rather than the sharper
$O_d(N\log(D+1))$ bound of the general fixed-bond theorem.
\end{theorem}

\begin{proof}
\subsection*{The exact support components}

View the support of $U$ as a bipartite graph, with row words
on one side and column words on the other.
If a row meets two columns, their inner product cannot vanish
through that one row alone. Since both columns have degree
at most two, they must share a second row.
Both rows and columns are now saturated. This is an isolated
complete $2\times2$ component. All remaining components are
singleton edges: a column incident to two rows would similarly
force a second column. Thus $U$ is a direct sum of unitary
$1\times1$ and $2\times2$ blocks after separate row and column
permutations.

In a $2\times2$ unitary block, the two diagonal moduli agree,
the two off-diagonal moduli agree, and the sum of their
squares is one. At least one of its two perfect matchings
therefore consists of entries of modulus at least $1/\sqrt2$.
Every row and column has such a dominant entry.

\subsection*{One common numerical graph}

Choose a positive dyadic $\tau$ with
$\varepsilon/200\le\tau\le\varepsilon/100$;
it requires $O(\log(1/\varepsilon)+1)$ bits.
Using a right-canonical MPS for the normalized Choi vector,
choose one deterministic fixed-point scalar calculation
$v(x,y)$ satisfying
\[
 |v(x,y)-U_{x,y}|\le\tau/10
       \qquad\text{for every }x,y.
\]
The canonical matrix norms are at most one. Their scalar
contraction estimate, together with the rescaling by
$d^{N/2}$, gives this accuracy using
\[
 B=O_d\bigl(N+\log(N+1)+\log(D+1)
                       +\log(1/\varepsilon)+1\bigr)
\]
bits and polynomially many Boolean operations.
Use this same calculation in all row and column queries.
Define a single exact Boolean edge predicate
\[
 (x,y)\in\widehat E\quad\Longleftrightarrow\quad
                   |v(x,y)|>2\tau.
\]
The comparison is made by squaring the real and imaginary
dyadic values and comparing integers. Thus it is an exact
predicate even when a value lies close to the threshold.
A true zero is never an edge of $\widehat E$.
Every retained edge has true modulus greater than $1.9\tau$,
and every dominant edge belongs to $\widehat E$.

The neighbor lists of this graph are computable with the
same polynomial precision.
For a fixed column, the normalized Choi mass of an output
prefix is
$d^{-N}\sum|U_{x,y}|^2$ over its extensions.
Conditioned right-canonical transfers are subunital and
contract the operator norm on Hermitian matrices.
Hermitian fixed-point rounding therefore computes these
masses with uniform error at most $\tau^2d^{-N}/100$,
using the stated $B$-bit precision.
Retain a prefix when its numerical mass exceeds
$\tau^2d^{-N}$. A retained prefix has positive true mass,
so there are at most two retained prefixes at each level.
Every edge of $\widehat E$ has modulus greater than
$1.9\tau$, hence all its prefixes have mass greater than
$3.61\tau^2d^{-N}$ and are retained.
At the final level, filter the candidate leaves by the
common predicate $\widehat E$.

For row queries use the corresponding canonical tensors
of $U^\dagger$ only for prefix-mass calculations.
Filter their candidate column words using the original
$v(x,y)$, not a separately rounded conjugate calculation.
Completeness of the searches then proves that both
orientations list exactly the same edges of $\widehat E$.
The number of prefix tests is $O_d(N)$ per query.
Each transfer and scalar calculation has polynomial
cost in $N,D,B$.

The graph $\widehat E$ is consequently an actual subgraph
of the exact support, containing a dominant perfect matching
in every genuine $2\times2$ block.
If no minor edge is retained, that block gives two singleton
components. If at least one minor edge is retained, its two
dominant matching edges connect all four vertices.
Thus every component of $\widehat E$ has equally many row
and column vertices, either one of each or two of each.
These are genuine components of one fixed graph; no
agreement between separately rounded reciprocal moduli
is assumed.

\subsection*{Two compatible address maps}

For a component, sort its column words as
$c_0$ or $(c_0,c_1)$, and its row words as
$r_0$ or $(r_0,r_1)$.
Its owner is the column word $c_0$.
Bounded-depth traversal of the graph computes these lists
from any of its row or column vertices.
Since a component has at most four vertices, this uses
only constantly many neighbor queries.
Define
\[
 g_{\rm in}(c_i)=(c_0,i),\qquad
 g_{\rm out}(r_i)=(c_0,i).
\]
They are bijections from the input and output basis word
sets, respectively, onto the same valid address set
\[
 \mathcal M=\{(c_0,i):
       c_0\text{ is a component owner and }
       0\le i<\text{its component size}\}.
\]
Their inverse functions on $\mathcal M$ are
$h_{\rm in}(c_0,i)=c_i$ and
$h_{\rm out}(c_0,i)=r_i$.
On invalid addresses extend these functions arbitrarily,
for instance with output zero. All four functions have
polynomial-size Boolean circuits.

Let $Y$ be the physical word register, let $O$ be an
additional $N$-qudit owner register, and let $I$ be one
qudit whose levels $0,1$ hold the component index.
A reversible encoder $E_{\rm in}$ computes
$g_{\rm in}(Y)$ into initialized $O,I$, reverses its
Boolean work, and subtracts $h_{\rm in}(O,I)$
componentwise modulo $d$ from $Y$.
It has the exact initialized action
\[
 E_{\rm in}|c_i,0,0\rangle=|0,c_0,i\rangle.
\]
Define $E_{\rm out}$ with $g_{\rm out},h_{\rm out}$:
\[
 E_{\rm out}|r_i,0,0\rangle=|0,c_0,i\rangle.
\]
Each subtraction is a reversible controlled addition,
even on invalid addresses. The encoders are therefore
full logical unitaries, not partial inverses supplied
as hypotheses. Their arithmetic work erases for all
logical inputs, since the inputs controlling each
Boolean computation are preserved while its result
is added to or subtracted from another register.

\subsection*{Exactly unitary approximations on each component}

For a two-row, two-column component, the actual matrix
\[
 B_{c_0}=(U_{r_i,c_j})_{0\le i,j<2}
\]
is an entire genuine $2\times2$ unitary block.
Compute its four entries to sufficient accuracy and
choose an exactly unitary approximation
$\widetilde B_{c_0}$ with
$\|\widetilde B_{c_0}-B_{c_0}\|\le\tau$.
Here is an explicit uniformly robust construction.
Write its first column as $(a,c)^T$, and put
$\zeta=\det B_{c_0}$. Since the matrix is unitary,
its second column is $\zeta(-\overline c,\overline a)^T$.
Choose phases $\alpha,\gamma,\beta$ and an angle $\theta$
such that
$a=e^{i\alpha}\cos\theta$,
$c=e^{i\gamma}\sin\theta$,
$\zeta=e^{i\beta}$, and $0\le\theta\le\pi/2$.
Then
\[
 B_{c_0}=
 \operatorname{diag}(e^{i\alpha},e^{i\gamma})
 \begin{pmatrix}\cos\theta&-\sin\theta\\
                 \sin\theta& \cos\theta\end{pmatrix}
 \operatorname{diag}(1,e^{i(\beta-\alpha-\gamma)}).
\]
This also specifies a unitary for arbitrary approximate
values of its four angles.

Compute the matrix entries more finely, with error
$c\tau^2$ for a sufficiently small fixed constant $c$.
Find $\theta$ from the ratio of the smaller and larger
first-column moduli: use its arctangent when $|a|\ge|c|$,
and $\pi/2$ minus its arctangent otherwise.
The larger approximate modulus is at least
$1/(2\sqrt2)$ in the approximate data.
Its arctangent is therefore uniformly conditioned.
For a component whose modulus estimate is at most
$\tau/100$, set its phase to zero: the corresponding
true modulus is $O(\tau)$, so this changes the
first column and its conjugate second column by
only $O(\tau)$. Otherwise compute its phase to
$c\tau$ accuracy. The largest Cartesian component
used for its phase division is at least a fixed
multiple of $\tau$, and the finer entry precision
makes this calculation stable.
The determinant has modulus one, so its phase
calculation has a denominator bounded away from zero.
Choose the clipping and angle constants small enough
that the total matrix error is at most $\tau$.
Square roots, division, and argument evaluation have
polynomial bit complexity by bisection and the
geometrically convergent half-angle arctangent series.

Store the four approximate angles as
$O(\log(1/\tau))$-bit strings. The last diagonal angle
is their exact fixed-point difference
$\widehat\beta-\widehat\alpha-\widehat\gamma$;
constant integer guard bits suffice.
Implement each phase and plane rotation by the
product of fixed-angle gates controlled by its
binary bits. Thus the implemented matrix is exactly
unitary, with the prescribed approximate scalar phase;
no table indexed by exponentially many component
owners is required.

For a singleton component use the phase of its retained
entry, approximated by an exactly unitary scalar.
The retained entry is dominant, so its modulus is
at least $1/\sqrt2$. Its phase evaluation is uniformly
conditioned.
An original singleton block then has phase error
at most $\tau$.
If an original $2\times2$ block was split into two
singletons, both omitted minor entries have modulus
at most $2.1\tau$. Write their common modulus as $s$.
Each retained entry has modulus $\sqrt{1-s^2}$.
Replacing it by its unit-modulus phase changes it by
$1-\sqrt{1-s^2}\le s^2$.
Including phase approximation, the resulting block
error is at most
$s+s^2+\tau\le4\tau$.

Compute the entries, component size, and angle bits
from the owner $c_0$, apply the resulting exactly
unitary operation to the component-index qudit, and
reverse this calculation. The owner register is
unchanged, so all arithmetic work erases exactly.
For a two-dimensional component the operation preserves
the levels $0,1$, acting as identity on other levels.
For a singleton it multiplies level zero by the
computed phase and leaves other levels arbitrary
unitarily, for example unchanged.
On invalid owner words choose identity.
Thus this controlled operation $R$ is a full logical
unitary and preserves the address subspace
$\mathcal M$ exactly.

The complete physical implementation is
\[
 C=E_{\rm out}^\dagger R E_{\rm in}.
\]
It has, on initialized auxiliary registers, the
exact action
\[
 |c_j,0,0\rangle
      \longmapsto
 \sum_i(\widetilde B_{c_0})_{ij}|r_i,0,0\rangle.
\]
The same holds for singleton components.
Consequently $CJ=J\widetilde U$ for an exactly unitary
$\widetilde U$. All owner, index, and arithmetic
registers return to zero.
The comparison with $U$ is blockwise on its genuine
support components. The global operator error is
the maximum block error, at most $4\tau<\varepsilon$;
there is no factor from the exponentially many
basis words.

All searches, address calculations, and angle
calculations have polynomial bit complexity in
$N,D,\log(1/\varepsilon)$. Retaining and reversing
their Boolean intermediate results gives polynomial
gate count and clean workspace. Routing through this
polynomial-size register adds only a polynomial factor.
\end{proof}

The support-two argument avoids polar synthesis of a
large rounded physical matrix by reducing it to
constant-dimensional blocks and preserving their
address subspace exactly. Such a reduction has not
been proved here for higher fixed sparsity. In
particular, the theorem does not resolve the joint
polynomial circuit question for general dense MPUs.

\begin{thebibliography}{3}
\bibitem{fixed-alphabet}
TNLean contributors,
\emph{Exact circuits for MPUs with a fixed physical-entry alphabet},
3 October 2026.
Local source:
\nolinkurl{docs/audits/2026-10-03_mpu_fixed_physical_alphabet.tex}.
\bibitem{permutation-circuits}
TNLean contributors,
\emph{Joint polynomial circuits for permutation and monomial unitaries},
3 October 2026.
Local source:
\nolinkurl{docs/audits/2026-10-03_mpu_joint_polynomial_permutations.tex}.
\bibitem{exact-mpu-circuits}
TNLean contributors,
\emph{Exact polynomial circuits for matrix product unitaries},
2 October 2026, Section~5.
Local source:
\nolinkurl{docs/audits/2026-10-02_mpu_rank_two_circuits.tex}.
\end{thebibliography}
\end{document}
